\uuid{Chzn}
\titre{Zéros de fonction}
\auteur{bodin ; hanani ; mzari}
\organisation{exo7}
\licence{CC BY-NC-SA 4.0}
\datecreate{2020-01-20}
\niveau{L1}
\module{Analyse}
\chapitre{Continuité, limite et étude de fonctions réelles}
\sousChapitre{Continuité : théorie}
\difficulte{2}
\theme{Théorèmes des valeurs intermédiaires, zéros de fonction}
\exo7qcmid{lille-1.continuite.022}

\contenu{
\begin{question}
\qclassification{123.01,123.02}
\qtags{motcle=zéros de fonction}
Soit $f : [0,1] \to \Rr$ une fonction continue avec $f(0)<0$ et $f(1)>0$.
Par dichotomie on construit deux suites $(a_n)$ et $(b_n)$, avec $a_0 = 0$ et $b_0 = 1$. Quelles sont les assertions vraies ?
\begin{answers}
    \bad{Si $f(\frac12)>0$ alors $a_1 = \frac12$ et $b_1 = \frac12$.}

    \good{$f$ s'annule sur $[a_n,b_n]$ (quel que soit $n\ge0$).}

    \bad{$(a_n)$ et $(b_n)$ sont des suites croissantes.}

    \bad{$a_n \to 0$ ou $b_n \to 0$.}
\end{answers}
\begin{explanations}
Par dichotomie on construit deux suites adjacentes. $(a_n)$ est croissante, $(b_n)$ est décroissante et $a_n \le b_n$. Ces deux suites convergent vers une valeur $c\in]a,b[$, telle que $f(c)=0$.
\end{explanations}
\end{question}
}
